\documentclass[12pt]{article}
\usepackage[a4paper,margin=2.5cm]{geometry}
\usepackage[T1]{fontenc}
\usepackage[utf8]{inputenc}
\usepackage{lmodern}
\usepackage{microtype}
\usepackage{needspace}
\usepackage{booktabs}
\usepackage{xcolor}
\usepackage{listings}
\usepackage[hidelinks]{hyperref}
\definecolor{codegray}{RGB}{245,245,245}
\lstset{basicstyle=\ttfamily\small,backgroundcolor=\color{codegray},frame=single,
  breaklines=true,showstringspaces=false,columns=fullflexible,keepspaces=true}
\setlength{\parindent}{0pt}
\setlength{\parskip}{0.55em}

\begin{document}
\pagenumbering{gobble}
\begin{titlepage}
\centering
{\large Ministerul Educa\c{t}iei \c{s}i Cercet\u{a}rii al Republicii Moldova\par}
\vspace{0.45cm}
{\large Universitatea Tehnic\u{a} a Moldovei\par}
\vspace{0.45cm}
{\large Facultatea Calculatoare, Informatic\u{a} \c{s}i Microelectronic\u{a}\par}
\vspace{0.45cm}
{\large Departamentul Ingineria Software \c{s}i Automatic\u{a}\par}
\vspace{3cm}
{\Large\bfseries Report on Laboratory Work No. 9\par}
\vspace{0.55cm}
{\LARGE\bfseries Watching System Calls with strace\par}
\vspace{0.35cm}
{\large Operating Systems\par}
\vfill
\begin{flushleft}
\textbf{Done by:} st. gr. FAF-241, Ca\c{t} Chiril\\[0.8cm]
\textbf{Checked by:} asist. univ. Rei\c{t}man Patricia
\end{flushleft}
\vfill
{Chi\c{s}in\u{a}u -- 2026\par}
\end{titlepage}
\pagenumbering{roman}
\tableofcontents
\clearpage
\pagenumbering{arabic}

\section{Aim, environment, and deliverables}
The aim is to observe the boundary between a user program and the Linux
kernel. I traced a real command, implemented a small file round trip using
\texttt{open}, \texttt{write}, \texttt{read}, and \texttt{close}, and counted
its write calls. Unlike the preceding xv6 exercise, this lab runs on Linux.

All compilation, execution, tracing, and automated tests ran on Ubuntu
24.04.5 LTS on an ARM64 Linux system, using GCC 13.3.0 and strace 6.8.
The compiler completed with no warnings under
\texttt{-std=c11 -Wall -Wextra -Werror}.

The supplied handout has two pages. Its overview names Part 3, ``Count one
kind of call'', but there are no detailed Part 3 instructions or Observe
questions in the supplied PDF. The counting experiment below is therefore
a supplementary interpretation of that heading, not a quotation of missing
instructions. Both Observe questions that actually appear in the handout
are answered explicitly in this report and in \texttt{lab9/report.md}.

\begin{lstlisting}
# On Ubuntu, compile the program:
gcc -std=c11 -Wall -Wextra -Werror -O0 -g \
  -o syscalls syscalls.c
./syscalls
strace -e trace=openat,write,read,close ./syscalls
\end{lstlisting}

The \texttt{code/} directory contains the C source and supporting
verification scripts. Recorded evidence includes raw traces, standard output, call
summaries, the build log, environment information, and verification results.
These are saved under \path{lab9/report_files/output/}.
The runner creates a fresh one-file listing directory and sets
\texttt{LC\_ALL=C}, \texttt{TZ=UTC}, and \texttt{umask 022} so the experiment
does not depend on unrelated host files. It records strace output with
\texttt{-o} separately from the program's standard output.

\clearpage
\section{Part 1: tracing a real command}
On Ubuntu, in the prepared directory, I ran:
\begin{lstlisting}
strace -f -e trace=openat,read,write,close -o ls.trace ls -l
strace -c -o ls.summary ls -l > /dev/null
\end{lstlisting}
The runner uses absolute output paths; the commands above abbreviate those
paths for readability. The following lines are excerpts of the actual
filtered trace. The PID prefix was 22 in this run:
\begin{lstlisting}[basicstyle=\ttfamily\footnotesize]
22 openat(AT_FDCWD, "/etc/ld.so.cache", O_RDONLY|O_CLOEXEC) = 3
22 close(3) = 0
22 openat(AT_FDCWD, ".", O_RDONLY|O_NONBLOCK|O_CLOEXEC|O_DIRECTORY) = 3
22 openat(AT_FDCWD, "/etc/passwd", O_RDONLY|O_CLOEXEC) = 4
22 read(4, "root:x:0:0:root:/root:/bin/bash\n"..., 4096) = 888
22 close(4) = 0
22 write(1, "total 4\n-rw-r--r-- 1 root root 2"..., 59) = 59
22 close(1) = 0
22 close(2) = 0
\end{lstlisting}

\subsection{What the calls mean}
The first calls load shared libraries and configuration. Opening
\texttt{"."} obtains a descriptor for the current directory. The long
listing also needs user and group information, which explains reads from
\texttt{/etc/passwd} and \texttt{/etc/group}. The listing itself is sent
with \texttt{write(1, ...)}. Descriptor 1 is standard output: on an
interactive terminal it leads to the screen, while in this captured run it
was redirected to \texttt{ls.stdout}.

The filter is not a complete account of directory traversal. Linux reads
directory entries with \texttt{getdents64}; the full summary contains two
such calls, but the handout's four-call filter hides them. Metadata calls,
memory mapping, and process startup also disappear from that filtered view.

\subsection{How many calls were observed?}
The filtered trace contained \textbf{38} calls: 10 \texttt{openat},
11 \texttt{read}, 16 \texttt{close}, and one \texttt{write}. The final
unfiltered \texttt{strace -c} run reported \textbf{113} calls in total,
including 12 errors. These include unsuccessful environment probes, not
necessarily failures of the listing command, which exited successfully.
An earlier run reported 114 total calls. Thus the useful rough answer is
about a hundred calls for this small listing, not a universal constant.
Architecture, libc, configuration, output destination, and directory
contents can change the count. Timing columns in the raw summary are not
used as performance measurements.

\clearpage
\section{Part 2: writing and reading through descriptors}
\subsection{Implementation and normal output}
The program follows the handout's sequence: open \texttt{note.txt} for
writing, send the message, close it, reopen it for reading, and display
what was read. The file itself is accessed using descriptor-based
\texttt{write} and \texttt{read}, not \texttt{fprintf} or \texttt{fread}.
The important operations are:
\begin{lstlisting}[language=C]
const char *msg = "hello, kernel\n";
int fd = open("note.txt", O_CREAT | O_WRONLY | O_TRUNC, 0644);
/* Check fd, then write the message, retrying interruptions
   and handling short writes. Close before reopening. */
fd = open("note.txt", O_RDONLY);
char buf[64];
n = read(fd, buf, sizeof(buf) - 1);
/* Check n before using it as an index. */
buf[n] = '\0';
printf("read got %zd bytes: %s", n, buf);
close(fd);
\end{lstlisting}
This is an abbreviated explanation, not the complete error-handling code.
The complete \texttt{syscalls.c} checks both opens, writes, read, and both
normal-path closes. It retries interrupted I/O and handles short writes.
In particular, it does not index \texttt{buf[-1]} after a failed read, a
risk in the unchecked handout example.

\begin{lstlisting}
$ gcc -std=c11 -Wall -Wextra -Werror -O0 -g \
    -o syscalls syscalls.c
$ ./syscalls
open returned descriptor 3
write sent 14 bytes
read got 14 bytes: hello, kernel
\end{lstlisting}

The message contains 14 bytes, including the newline but excluding C's
terminating null byte. \texttt{O\_TRUNC} removes previous contents of an
existing file; \texttt{O\_CREAT} creates it if necessary. The requested mode
is 0644, with actual permissions also subject to the process umask.
\texttt{read} requests at most 63 bytes, leaving space for the added null
terminator, and returns the 14 bytes available in this file. A successful
read in general can return fewer bytes than requested; this example is a
single read of the tiny local file, not a general-purpose file copier.

\clearpage
\subsection{The observed kernel calls}
\lstinputlisting[basicstyle=\ttfamily\footnotesize]{output/syscalls.trace}
The first five lines concern the loader and libc. The next six lines
correspond directly to the file round trip. \texttt{AT\_FDCWD} means a
relative pathname is resolved against the current working directory.
The C library's \texttt{open()} wrapper uses the Linux \texttt{openat}
system call in this environment. These are libc wrappers around kernel
services, not handwritten assembly or calls to the \texttt{syscall()}
function.

The first file descriptor is 3 because 0, 1, and 2 were already inherited
as standard input, standard output, and standard error. After
\texttt{close(3)}, that number is free, so the second open reuses 3.
A descriptor is a per-process integer handle, not a permanent file ID.

\clearpage
\section{Part 3: counting one kind of call (supplement)}
I chose \texttt{write} to compare a C-level print with actual kernel calls:
\begin{lstlisting}
strace -c -e trace=write -o write.summary ./syscalls
strace -e trace=write -o write.trace ./syscalls
\end{lstlisting}
\lstinputlisting[basicstyle=\ttfamily\footnotesize]{output/write.trace}

The independent summary also reports two \texttt{write} calls:
\lstinputlisting[basicstyle=\ttfamily\footnotesize]{output/write.summary}

One call sends the 14-byte message to descriptor 3. The other sends all
80 bytes of diagnostic output to descriptor 1. Although the source calls
\texttt{printf} three times, standard output was redirected to a regular
file during evidence capture and libc buffered those messages until exit.
The file write therefore appears before the diagnostic write in the trace,
even though the first \texttt{printf} appears earlier in the source.
On a terminal, stdout is normally line buffered, so the number and ordering
of stdout writes can differ. Counting \texttt{printf} statements is not
how to count \texttt{write} system calls.

\section{Handout questions and answers}
\begin{enumerate}
\Needspace{8\baselineskip}
\item \textbf{Question:} Which system call does \texttt{ls} use to print to the screen? To which descriptor does it write?
\par\textbf{Answer:} It uses \texttt{write} to descriptor 1, standard output. The recorded trace shows a 59-byte \texttt{write(1, ...)}. Here stdout was captured in a file rather than displayed on a terminal; the descriptor's role is unchanged.

\Needspace{8\baselineskip}
\item \textbf{Question:} The descriptor returned by the first \texttt{openat} is 3, not 0, 1 or 2. Why are those three already taken?
\par\textbf{Answer:} The process inherits descriptor 0 for standard input, 1 for standard output, and 2 for standard error. With those open, 3 is the lowest available descriptor when \texttt{note.txt} is opened. Library startup briefly used 3 and closed it before \texttt{main} reached the file open.
\end{enumerate}

\clearpage
\section{Verification and conclusion}
The retained verifier checked the program's exact output in normal, traced,
and call-counting runs, and compared \texttt{note.txt} with the expected
14 bytes. It checked all six file calls in order, matched their return
values, and independently compared the two-write trace with the summary.
It also tested these cases in a temporary directory:
\begin{itemize}
\item An existing longer file is truncated and replaced by exactly the message.
\item A directory named \texttt{note.txt} causes the write-open to fail; the program reports the error and exits with status 1.
\item A \texttt{note.txt} symlink to \texttt{/dev/full} forces a write failure; the program reports it and exits with status 1.
\end{itemize}
All \textbf{15 assertions} passed on the final complete rerun.
The verification evidence is in \texttt{output/verification.txt}; the
complete shell command log and environment are in \texttt{output/session.txt}.
The tests cover normal behavior, truncation, and representative open/write
errors, but do not claim exhaustive fault injection for every read or close
failure.

The experiment connects source-level operations to observed Linux calls.
The kernel returns small descriptor numbers that can be reused after
closing; calls return actual byte counts rather than the requested amount
by definition. A simple command performs many startup and metadata calls
in addition to its visible output. Finally, stdio buffering explains why
several C print operations can become one kernel write, while the explicit
file write appears immediately in the trace.

\end{document}
